% !TEX root = A0.MiTFG.tex

\chapter{Introducción}

\section{Contexto y motivación}
\apartadoPendiente

\subsection{Crecimiento del \textit{streaming} OTT}
\apartadoPendiente

\subsection{Importancia del DRM}
\apartadoPendiente

\subsection{Widevine como estándar dominante}
\apartadoPendiente

\subsection{Limitaciones de asumir que tener DRM implica seguridad}
\apartadoPendiente

\section{Problema de investigación}
\apartadoPendiente

\subsection{Qué protege realmente Widevine}
\apartadoPendiente

\subsection{Qué no protege Widevine}
\apartadoPendiente

\subsection{Amenazas persistentes en ecosistemas OTT}
\apartadoPendiente

\subsubsection{CDN \textit{leeching}}
\apartadoPendiente

\subsubsection{Reutilización de tokens}
\apartadoPendiente

\subsubsection{\textit{Credential sharing}}
\apartadoPendiente

\subsubsection{Accesos concurrentes}
\apartadoPendiente

\subsubsection{\textit{Bypass} de controles perimetrales}
\apartadoPendiente

\section{Hipótesis de trabajo}
\apartadoPendiente

\section{Objetivos}
\apartadoPendiente

\subsection{Objetivo general}

El objetivo principal de este TFM es diseñar, implementar y auditar un entorno experimental de \textit{streaming} de vídeo en directo que integre el sistema de protección Widevine DRM, con el fin de evaluar sus vulnerabilidades inherentes y proponer una arquitectura de defensa en profundidad mediante el desarrollo de capas adicionales de autenticación, autorización dinámica y control de concurrencia.

\subsection{Objetivos específicos}

\begin{enumerate}
    \item Desplegar una infraestructura completa de distribución OTT (\textit{Over-The-Top}) mediante contenedores, simulando la cadena completa de vídeo, incluyendo el empaquetado con cifrado Common Encryption (CENC), un servidor de origen y una red de distribución de contenidos (CDN).
    \item Realizar una auditoría de seguridad del ecosistema DRM Widevine, analizando el flujo de intercambio de licencias entre el cliente, el Content Decryption Module (CDM) y el servidor de licencias, e identificando vectores de ataque críticos.
    \item Desarrollar un proxy para gestionar las licencias del contenido con lógica de autorización dinámica, interceptando las peticiones de licencia y validando en tiempo real la identidad del usuario mediante tokens JWT, su ubicación geográfica y el estado de su suscripción antes de permitir la comunicación con el servidor de licencias.
    \item Implementar mecanismos adicionales de seguridad que limiten el abuso, como controles de concurrencia y detección de anomalías capaces de bloquear automáticamente sesiones o direcciones IP al detectar patrones de uso compartido de cuentas o intentos de acceso masivo.
\end{enumerate}

\section{Metodología}
\apartadoPendiente

\subsection{Metodología incremental}
\apartadoPendiente

\subsection{Laboratorio experimental}
\apartadoPendiente

\subsection{Auditoría técnica}
\apartadoPendiente

\subsection{Diseño iterativo de defensas}
\apartadoPendiente

\section{Estructura de la memoria}
\apartadoPendiente
